\section{Alternating character products}\label{sec:products}

For the rest of the argument, fix an odd integer \(u>1\), put
\(P=C_{u^2}\), and let \(I\) be the nonempty set of primes dividing
\(u\). Decompose \(P\) into its cyclic primary components. For each
\(p\in I\), choose a generator \(X_p\) of the \(p\)-component and a
primitive \(p\)-th root \(\rho_p\). These choices remain fixed.
For \(S\subseteq I\), let \(\chi_S\) be the character determined by
\[
 \chi_S(X_p)=
 \begin{cases}
  \rho_p,&p\in S,\\
  1,&p\notin S,
 \end{cases}
 \qquad
 x_S=\sum_{z\in P}x_z\chi_S(z)
 \quad\text{for }x=\sum_{z\in P}x_z z\in\Z[i][P].
\]
All these values lie in the same ring
\begin{equation}\label{eq:alternating-notation}
 B=\Z[i,\{\rho_p:p\in I\}],\qquad
 \eps_S=(-1)^{|S|}.
\end{equation}

If \(x\in\Z[i][P]\) satisfies \(xx^*=u^2\), character evaluation gives
\begin{equation}\label{eq:character-norm}
 x_S\overline{x_S}=u^2.
\end{equation}
In particular every \(x_S\) is nonzero, so we may form
\begin{equation}\label{eq:delta}
 \Delta(x)=\prod_{S\subseteq I}x_S^{\eps_S}\in\Frac(B).
\end{equation}
This product combines comparisons between pairs of characters. When
\(I=\{p,q\}\), for example,
\[
 \begin{aligned}
 \Delta(x)
  &=\frac{x_{\varnothing}x_{\{p,q\}}}
          {x_{\{p\}}x_{\{q\}}}\\
  &=\frac{x_{\varnothing}/x_{\{p\}}}
          {x_{\{q\}}/x_{\{p,q\}}}
   =\frac{x_{\varnothing}/x_{\{q\}}}
          {x_{\{p\}}/x_{\{p,q\}}}.
 \end{aligned}
\]
The two expressions on the second line pair the same four values in
the \(p\)- and \(q\)-directions. A comparison in one direction controls
primes above that prime; the same product admits a pairing in every
direction. We will use these pairings to make \(\Delta(x)\) integral.
The norm identity and Kronecker's criterion will then give finite order,
and the local residue-one conditions will force that order to be odd.

\subsection{Comparison of two character values}

For \(p\in I\), write \(p^e\) for the order of the \(p\)-component
of \(P\). Then
\(e=v_p(u^2)\): the scalar norm has exactly the same \(p\)-valuation
as the exponent in the component order. This is the case treated
by the following local statement. It applies to two separate factors,
and its proof uses only the sum of their valuations.

\begin{lemma}[Character comparison]\label{lem:comparison}
Let \(A\) be as in Lemma~\ref{lem:local-ring}. Suppose that \(e\geq1\)
and
\[
 F,K\in A[C_{p^e}],\qquad FK=p^e b,\qquad b\in A^\times.
\]
For a primitive \(p\)-th root \(\rho\), both ratios
\[
 \frac{F(\rho)}{F(1)},\qquad \frac{K(\rho)}{K(1)}
\]
are units in \(A[\rho]\) with residue \(1\) at its maximal ideal.
\end{lemma}

\begin{proof}
All the values in the denominators are nonzero, since the product
of the two corresponding values is the nonzero scalar \(p^e b\).
The ratios therefore begin as elements of \(\Frac(A[\rho])\).

For \(e=1\), the nonnegative integral valuations of \(F(1)\) and
\(K(1)\) sum to \(1\), so one is a unit of \(A\). Evaluation at
\(\rho\) is congruent to evaluation at \(1\) modulo \(\rho-1\).
For that factor, the ratio is therefore a unit of residue \(1\).
The two ratios multiply to \(1\), proving the same assertion for
the other factor.

Suppose \(e\geq2\). Let
\(q:A[C_{p^e}]\to A[C_{p^{e-1}}]\) be the projection sending the
chosen generator to the chosen generator. We will show that the
projection of one factor is divisible by \(p\) in the smaller group
ring. Evaluate at a primitive \(p^e\)-th root \(\xi\), and put
\(D=(p-1)p^{e-1}\) as in Lemma~\ref{lem:local-ring}. The two
nonnegative valuations sum to \(e\), so at least one is at least \(1\).
Interchange the factors if necessary so that
\(F(\xi)\in pA[\xi]\).

Choose a polynomial representative of \(F\) and divide by the monic
cyclotomic polynomial:
\[
 F=Q\Phi_{p^e}+R,\qquad Q,R\in A[X],\quad \deg R<D.
\]
Since \(R(\xi)=F(\xi)\), the coefficient statement in
Lemma~\ref{lem:local-ring} makes every coefficient of \(R\) divisible
by \(p\). Moreover \(q(\Phi_{p^e})=p\). Thus
\(q(F)\in pA[C_{p^{e-1}}]\), as required. Define
\[
 F'=q(F)/p,\qquad K'=q(K).
\]
Then \(F'K'=p^{e-1}b\). Cancellation of \(p\) is valid here because
the target group ring is a free \(A\)-module. The evaluations at
\(1\) and \(\rho\) factor through \(q\), and scalar division cancels
from their ratio. Thus both original comparison ratios are unchanged.
Induction proves the assertion, whether the same factor or the other
factor is divided at the next step.
\end{proof}

\begin{remark}\label{rem:exponent}
The scalar valuation in Lemma~\ref{lem:comparison} cannot uniformly
be allowed to exceed the exponent in the cyclic group order. For an
odd prime \(p\) and a primitive \(p\)-th root \(\rho\), put
\(J=1+X+\cdots+X^{p-1}\) in \(\Z[C_p]\) and \(F=p-2J\).
Since \(J^2=pJ\) and \(J^*=J\), one has \(FF^*=p^2\),
but \(F(\rho)/F(1)=-1\), whose residue is not \(1\).
Here the group has order \(p\), while the scalar has \(p\)-valuation
\(2\). In the case used above and below, both exponents start at
\(e\), and every projection is accompanied by division of one factor
by \(p\), reducing both exponents by one.
\end{remark}

\subsection{From local units to roots of unity}

The comparison supplies control at the primes dividing \(u\).
The norm identity supplies both the remaining local control and the
complex absolute values. We use two elementary facts to turn these
properties into an order restriction. The first is Kronecker's
criterion \cite[Statement~I]{Kronecker1857}.

\Needspace{7\baselineskip}
\begin{lemma}\label{lem:torsion}
\leavevmode
\begin{enumerate}[label=\textup{(\roman*)}]
\item Let \(B\subseteq\C\) be generated over \(\Z\) by finitely many
roots of unity. If \(z\in B\) satisfies \(|\sigma(z)|=1\) for every
unital ring homomorphism \(\sigma:B\to\C\), then \(z\) is a root of unity.
\item In a local domain with residue characteristic \(p\), a root of
unity with residue \(1\) has \(p\)-power order. In particular, if its
order is coprime to \(p\), it equals \(1\).
\end{enumerate}
\end{lemma}

\begin{proof}
For (i), let \(d\) be the finite rank of the free additive group of
\(B\), as in the proof of Lemma~\ref{lem:local-ring}. Multiplication by an element
of \(B\) is thus an integer matrix. On complexifying this
representation, the matrices for the generating roots of unity
commute and have finite order. Each is diagonalizable, since its
minimal polynomial divides some \(X^m-1\), which has distinct complex
roots. They are simultaneously diagonalizable: the eigenspaces of one
are invariant under the others, and one repeats the argument on those
eigenspaces.

On a common eigenvector, multiplication by any \(b\in B\) acts as a
scalar \(\sigma(b)\). This is a well-defined unital ring homomorphism
because it comes from the multiplication representation, so all
relations between the root generators are respected. The multiplication
matrix of \(z\) is diagonal in the same basis, with eigenvalues
\(\lambda_1,\ldots,\lambda_d\) of modulus one. For each positive
integer \(a\), its \(a\)-th power has an integer characteristic
polynomial. The coefficient of degree \(d-j\) has absolute value at
most \(\binom dj\). Only finitely many such polynomials can occur,
so the union of their sets of roots is finite. Each sequence
\(\lambda_j^a\) therefore takes only finitely many values. Since
\(\lambda_j\ne0\), two equal powers imply that \(\lambda_j\) has
finite order. A common multiple of these
orders makes the multiplication matrix the identity; applying it to
\(1\in B\) proves that a positive power of \(z\) is \(1\).

For (ii), write the order as \(p^a m\) with \(p\nmid m\). If \(m>1\),
raising the root to the \(p^a\)-th power gives an element \(\zeta\)
of order \(m\) and residue \(1\). In a domain,
\((\zeta-1)(1+\zeta+\cdots+\zeta^{m-1})=0\) then implies
\[
 1+\zeta+\cdots+\zeta^{m-1}=0.
\]
Its reduction says \(m=0\) in the residue field, a contradiction.
Thus \(m=1\).
\end{proof}

\begin{proposition}\label{prop:alternating}
Suppose \(x\in\Z[i][P]\) satisfies \(xx^*=u^2\). Then the alternating
product \(\Delta(x)\) defined in \eqref{eq:delta} belongs to \(B\) and
is a root of unity. For each \(p\in I\), its order is a power of \(p\).
In particular, its order is odd; if \(u\) has at least two distinct
prime divisors, then \(\Delta(x)=1\).
\end{proposition}

\begin{proof}
Fix \(p\in I\) and a maximal ideal \(\m\) of \(B\) containing \(p\).
Let
\[
 A_0=\Z[i,\{\rho_q:q\in I\setminus\{p\}\}],\qquad
 \m_0=\m\cap A_0,\qquad A=(A_0)_{\m_0}.
\]
Put
\[
 \eta=i\prod_{q\in I\setminus\{p\}}\rho_q,\qquad
 L=4\prod_{q\in I\setminus\{p\}}q,
\]
with empty products equal to \(1\). The root orders \(4\) and
\(q\in I\setminus\{p\}\) are pairwise coprime, so \(\eta\) is a
primitive \(L\)-th root, \(p\nmid L\), and each original root is a
power of \(\eta\). Thus \(A_0=\Z[\eta]\) has the form
required by Lemma~\ref{lem:local-ring}. The ideal \(\m_0\) is maximal:
it is prime and contains \(p\), and \(A_0/pA_0\) is a product of fields
as in that lemma.

Let \(p^e\) be the order of the \(p\)-component, so \(e=v_p(u^2)\).
For \(p\notin S\), evaluate every other component in \(xx^*=u^2\)
according to \(\chi_S\), leaving the \(p\)-generator as a variable.
This ring homomorphism produces two factors \(F,K\in A_0[C_{p^e}]\)
with
\[
 FK=p^e b,\qquad b=u^2/p^e\in A^\times.
\]
Both factors already have coefficients in \(A_0\): in the image of
\(x^*\), Gaussian coefficients are conjugated and the evaluated
roots are inverted, and these operations preserve \(A_0\).
After localization, Lemma~\ref{lem:comparison} therefore applies
to these two factors. It gives
\begin{equation}\label{eq:edge-comparison}
 \frac{x_{S\cup\{p\}}}{x_S}\in A[\rho_p]^\times,
 \qquad
 \frac{x_{S\cup\{p\}}}{x_S}\equiv1\pmod{\rho_p-1}.
\end{equation}

Every denominator inverted in \(A\) lies outside \(\m\), so the
natural inclusion of \(A[\rho_p]\) in \(B_{\m}\) is defined.
In the residue field of \(B_{\m}\), characteristic \(p\) gives
\((\rho_p-1)^p=\rho_p^p-1=0\). Hence \(\rho_p-1\) has residue zero,
and the units in \eqref{eq:edge-comparison} retain residue \(1\)
in \(B_{\m}\). Pairing \(S\) with \(S\cup\{p\}\) yields the identity
\begin{equation}\label{eq:pairing}
 \Delta(x)=\prod_{\substack{S\subseteq I\\p\notin S}}
 \left(\frac{x_S}{x_{S\cup\{p\}}}\right)^{\eps_S}.
\end{equation}
Thus \(\Delta(x)\) is a unit of residue \(1\) in \(B_{\m}\).
The left side of \eqref{eq:pairing} is the same element when the
pairing direction changes. Applying the argument for every \(p\mid u\)
and every maximal ideal over \(p\) gives this local conclusion at all
of those ideals.

At a maximal ideal containing none of the prime divisors of \(u\),
the scalar \(u^2\) is a unit. Equation~\eqref{eq:character-norm}
then makes every \(x_S\) a unit there. Consequently \(\Delta(x)\)
belongs to \(B_{\m}\) for every maximal ideal \(\m\) of \(B\).
The intersection of these localizations inside \(\Frac(B)\) is \(B\).
Indeed, if \(z\notin B\), the ideal
\[
 D_z=\{b\in B:bz\in B\}
\]
is proper. A maximal ideal \(\m\) containing it cannot have
\(z\in B_{\m}\): such membership would give \(z=a/s\) with
\(s\notin\m\), whereas \(sz=a\in B\) puts \(s\) in \(D_z\).
This proves \(\Delta(x)\in B\).

Every unital homomorphism \(\sigma:B\to\C\) sends the root generators to
roots of unity and commutes with complex conjugation on their integer
polynomials. Applying it to \eqref{eq:character-norm} gives
\(|\sigma(x_S)|=u\). In particular these images are nonzero.
Let \(P_0\) and \(Q_0\) be the products of the \(x_S\) with
\(\eps_S=1\) and \(\eps_S=-1\), respectively. Then
\(\Delta(x)=P_0/Q_0\), and the preceding nonvanishing gives
\(\sigma(Q_0)\ne0\). Applying \(\sigma\) to
\(Q_0\Delta(x)=P_0\) in \(B\) therefore evaluates the fraction
in \eqref{eq:delta}.
Since \(I\ne\varnothing\),
\[
 \sum_{S\subseteq I}\eps_S=(1-1)^{|I|}=0.
\]
It follows that \(|\sigma(\Delta(x))|=1\).
Lemma~\ref{lem:torsion}(i) makes \(\Delta(x)\) a root of unity.

For each \(p\in I\), there is a maximal ideal of \(B\) containing
\(p\): the additive group of \(B\) is free of positive finite rank,
so \(B/pB\ne0\). At that ideal \(\Delta(x)\) has residue \(1\), so
Lemma~\ref{lem:torsion}(ii) makes its order a power of \(p\).
The assertions about odd order and two distinct prime divisors follow.
\end{proof}

The single-prime conclusion can be nontrivial. If \(u=p^a\), where
\(p\) is an odd prime and \(a\geq1\), and \(X\) is the chosen generator
of \(C_{p^{2a}}\), then \(x=p^aX\) satisfies \(xx^*=u^2\), while
\(\Delta(x)=\rho_p^{-1}\) has order \(p\). The final argument will
therefore use odd torsion, without requiring each alternating product
to equal \(1\).
